\chapter{Resultados Preliminares y Diseño Experimental}

Al encontrarse esta investigación en una fase de desarrollo arquitectónico e integración de modelos, no se dispone aún de una métrica global de reconstrucción a escala distrital. Sin embargo, en este capítulo se presentan las \textit{Pruebas de Concepto} (PoC) iniciales que validan la viabilidad técnica de los módulos individuales propuestos. Asimismo, se define el diseño experimental riguroso y las métricas con las cuales se evaluará el sistema final contra el estado del arte.

\section{Pruebas de Concepto (Validación Modular)}

Para garantizar que el enfoque \textit{top-down} es aplicable a la arquitectura limeña y entornos no estructurados, se ejecutaron pruebas de aislamiento sobre los modelos neuronales que fundamentan las fases de abstracción.

\subsection{Filtrado de Elementos Dinámicos en Entornos Locales}
La primera hipótesis del sistema exige que el ruido dinámico pueda ser aislado eficientemente antes del rastreo (\textit{tracking}). Las pruebas preliminares utilizando modelos de segmentación instantánea (como YOLO) sobre metraje real capturado confirmaron la capacidad de inferir, en tiempo real, las máscaras delimitadoras de vehículos y peatones. Tal como se expuso en el Capítulo IV, este enmascaramiento resulta lo suficientemente robusto para garantizar que los cálculos geométricos de triangulación se alimenten exclusivamente de píxeles estáticos arquitectónicos.

\subsection{Inferencia de Wireframes con Co-PLNet}
El corazón de la propuesta geométrica depende de la identificación de límites estructurales asumiendo restricciones \textit{Manhattan-world}. Las pruebas iniciales ejecutadas con Co-PLNet sobre edificios institucionales (campus UTEC) y edificaciones residenciales (distrito de Barranco) arrojaron resultados muy favorables. El modelo demostró no solo ser capaz de detectar puntos de interés robustos, sino también de interconectarlos mediante tensores de líneas, proveyendo así la base de un grafo bidimensional coherente incluso en fachadas con alta repetición de texturas (ej. ventanas uniformes) donde los descriptores clásicos como SIFT suelen fallar por ambigüedad.

\section{Diseño Experimental y Métricas de Evaluación Futuras}

Para la culminación del proyecto, la arquitectura unificada será sometida a una evaluación cuantitativa exhaustiva, comparando su rendimiento frente al \textit{pipeline} estándar de la industria dominado por COLMAP (SfM + MVS).

\subsection{Métricas de Precisión Geométrica}
La validación de las poses estimadas de la cámara en el grafo $SE(3)$ se realizará midiendo el \textbf{Error Cuadrático Medio (RMSE)} de la Trayectoria Absoluta (ATE) frente a un \textit{Ground Truth} recolectado mediante telemetría GPS o LIDAR de precisión.
\begin{equation}
    \text{RMSE}_{\text{ATE}} = \sqrt{ \frac{1}{N} \sum_{i=1}^{N} \left\| \mathbf{t}_i - \mathbf{t}_i^* \right\|^2 }
\end{equation}
donde $\mathbf{t}_i$ representa la traslación estimada y $\mathbf{t}_i^*$ la trayectoria real.

\subsection{Eficiencia Computacional y Escalabilidad}
La principal justificación de este proyecto radica en la ruptura del cuello de botella temporal. Por tanto, se documentará meticulosamente:
\begin{itemize}
    \item \textbf{Latencia de Procesamiento:} Tiempo medio requerido (en segundos por fotograma) para pasar de imágenes RGB a la malla poligonal estructurada, comparado con el tiempo de densificación píxel a píxel del MVS.
    \item \textbf{Huella de Memoria (Reducción Paramétrica):} Contraste del peso de los archivos exportados. Se medirá en Megabytes (MB) el modelo de planos matemáticos en contraposición al modelo denso de millones de puntos requerido para el mismo sector urbano.
\end{itemize}

El cumplimiento de estas métricas certificará estadísticamente la superioridad computacional del enfoque \textit{MVS-free} para tareas de modelado a escala masiva.